\section{地方档案留下的南宋：田、船、庙与城}

\subsection{从一块田的界址开始}

南宋地方社会最不容易被宏大年表照亮的，往往是田界怎样确认、欠债怎样延期、谁为一次交易作保。契约里的四至、邻人、价钱、期限和见证，并非琐碎到可以略去的细节。它们告诉我们，田宅能够进入市场，是因为当事人把地形、亲属、邻里和地方承认写进了一份可保存的文本。战乱之后的置田、典卖或租佃尤其如此：一个迁入家庭要取得的不只是几亩地，还要取得在新社区中被承认的界址、担保与祭祀关系。

这类文书不能替我们描述整个江南。保存下来的往往来自书写、保管条件较好的地区，也较多记录有财产、有亲属网络或发生争议的人。契约写下的是签约时希望成立的关系，并不能保证收成、还款和占有后来都照原约进行。但它让“土地兼并”“社会流动”之类的概括有机会回到可检验的尺度：每一项权利都要由人、物、时间和地点组成，不能只由一个制度名词解释。

官府的税籍与私人契约又不是两套毫无关联的世界。田地被买卖、析分或荒废时，赋役归属、邻界、户籍与地方催征都可能成为问题；但国家分类并不自动吞没家庭的做法。有人会借宗族、保人、寺院或市镇的信用维持交易，有人则因无文书、无担保或迁徙而更难被稳定地登记。材料无法为所有人给出相同清晰度，正说明所谓地方秩序是不断被书写、认可和争执的过程。

\subsection{借贷把农时、市场和亲属连在一起}

借贷文书中的利息、期限和抵押，常使读者想立刻判断债主是否剥削、债户是否破产。更需要追问的，是借款发生在哪一个农时、以米钱还是布帛计值、担保者是谁、债务是否与婚姻、迁居或纳税相连。对缺少现钱的家庭，借贷可能是度过青黄不接、购置生产工具或支付丧葬的办法；对债权人和中介，它也可能是扩展影响的机会。两种情形可以同时存在，不能由一份格式化文书替全部关系下判决。

纸币和市场价格会改变这种关系，却不会把它们化为纯粹的货币问题。会子能进入支付与债务的某些环节，但不同地点的接受程度、折换条件和信息渠道不一。军费与临时征收增加时，家庭的实物、劳力和信用更可能一起被动员。会要中的币制与财政条目说明朝廷怎样理解这种压力，契约则偶尔显示某种压力怎样落到一处交易；二者之间仍隔着大量未保存的议价、拖延、互助与失败。

\subsection{港口的仓库、寺观和翻译}

港口不是一条从中国通向“海外”的单线门槛。泉州、明州和广州等地的航运，需要内河船、码头、仓储、牙人、翻译、地方官和能应对风季的水手。海外商人、侨居者和本地商户进入同一个城市，并不表示他们拥有同一种法律位置或同一种安全。市舶机关可以征税、验货和发给许可，仍无法亲自完成每一次装卸、借贷、雇工和信息交换。

《诸蕃志》及相近的地理贸易记述，给南宋读者保留了港口所见或所闻的海外物产、航路和政权。它们应当被当作作者在特定信息网络中的写作，而不是一张精确的航海图。作者亲见与转述、旧闻与当时情形必须分开；地名和货名的对应也常有不确定性。与港口遗址、沉船和海外出土陶瓷并读，能够让某些路线和货类更可信，却不能让我们计算所有贸易或把器物的发现地点等同于国家控制范围。

港口寺观在这种流动中占有特殊位置。它们可为商人、船员和本地施主提供祭祀、捐施、住宿、书写和记忆的场所；题记常把一次航海、一次修建或一笔施入固定下来。功德碑并不是运输账册，沉船也不是关税档案，但两者共同提示：海贸的社会基础不仅有官署和商号，还有能够跨越季节、语言和短暂停泊维持信任的地方组织。

\subsection{窑火、船木与不被计入的技术}

外销瓷、漆器和纺织品的光彩，容易遮住生产过程里的燃料、泥料、木材、运输和工匠。窑址的分布、匣钵和残器可以让考古学家辨认某些制作方式、产品类别与时间层次；它们不能直接告诉我们作坊的雇佣形式、所有劳动者的籍贯或利润如何分配。生产者在城镇与乡村之间流动，可能同时受市场订单、地方税役、灾害和军需影响。

造船也是如此。南宋能够使用江海运输与水军，依赖港湾、造船场、木材、铁器、绳索、熟练工和船员经验。技术并不是脱离社会的优势；它要由持续的供给和组织维持。战争加紧时，军方征用船只或工匠可以解决一时的防御需要，也可能挤压商运和地方劳力。对这些变化，官府记录较容易留下征发和制度，劳动者的议价与流动却较少留下完整的声音。

\subsection{城门以内的夜市与城门以外的供给}

城市消费在南宋笔记中留下了鲜明图像：夜市、酒楼、戏乐、节日和各色店铺。不过，城市并不是自身生产这些景象。柴薪、蔬菜、米粮、鱼盐、手工品和服务都要越过城门、桥梁、河埠与市场进入；管理者关心的仓储、消防、治安、度量和税收，正说明这种供给容易受阻。读笔记时不应把一场节日写成全城都享有的安稳，也不应因其文学性而忽视城市确实积累了复杂的商业和服务网络。

临安的繁华和前线的压力在同一财政体系中相遇。官员、军人和随从的消费提高了城市需求，江南的市场又提供税源与信用；一旦战线逼近、运输受阻或征发加重，原本连接城市和腹地的渠道便可能成为紧张点。城市的适应能力来自多种经营与水路联系，不来自对外部风险的免疫。市场仍在交易，并不意味着每个家庭都有余力承担价格和税役的变化。

\subsection{寺院经济的边界}

寺院拥有田地、接受施入、组织法会或保存经卷，常被概括成“寺院经济”。这个说法有用，但必须回到具体权利。施田是何人何时以何种名义给予，租佃和管理如何进行，地方官是否承认某种免役或保护，皆会改变一座寺院的实际资源。碑刻会突出布施者的虔诚和寺院的功德，官府材料会突出名籍、度牒和限制；两者不必相互否定，只是各自在不同目的下书写同一空间。

宗教组织能在迁徙与灾害时提供接待、安葬和互助的渠道，却不应因此被浪漫化为人人可及的安全网。寺院有自身的资源与等级，也要同地方权力、家族和官府交涉。国家对僧道的分类能够影响资格与税役，无法穷尽个人信仰、地方习俗或跨地域结社。保留这个差别，才能避免把一纸禁令写成信仰已经消失，或把一通造像记写成社会完全脱离行政。

\subsection{从地方材料看到末年的连续性}

襄樊失守、临安降元和崖山覆亡改变了政权的年号和官府，地方的田、船、庙、债务与亲属关系却不可能随年号同时停止。旧契约可能仍要被解释，港口仍要处理货物和人员，寺院仍要面对施主与新政权。元初材料所记录的接收、重新登记和任命，正说明新政府也必须借助已有的道路、财产关系和地方中介。它们不允许我们假设转变毫无阻碍，更不允许把后来被登记的人都看作最初便主动接受。

这些局部证据不能替代政权年表，却能纠正年表容易造成的错觉。南宋末年并非只有一条从败战走向覆亡的道路；在不同城镇、河口、乡村和寺院，人们处理的是不同的粮食、亲属、航路和文书问题。战争终结一个朝廷的能力，常常先以缓慢、分散的方式改变这些问题。把它们放回战争、财政和海贸的连续叙事，才能让“地方社会”不只是大事之后的附录。
